\title{Parameter-Efficient Orthogonal Finetuning via Factorization}

\begin{abstract}
With the widespread adoption of large foundation models, retraining these architectures from scratch has become increasingly infeasible due to their immense computational requirements. Consequently, there is a pressing need for methods that adapt such models to downstream tasks with minimal additional training. This study investigates Orthogonal Finetuning (OFT), a principled approach for adapting foundation models to new tasks. While OFT demonstrates promising generalization capabilities, it often necessitates a substantial number of trainable parameters owing to the inherently high dimensionality of orthogonal matrices. To overcome this, we analyze OFT through the lens of information transmission, deriving key criteria for achieving greater parameter efficiency. Drawing inspiration from the efficiency of the Cooley-Tukey fast Fourier transform algorithm in information transmission, we develop a compact orthogonal parameterization utilizing butterfly structures. Integrating this parameterization within OFT, we introduce Orthogonal Butterfly (BOFT), a novel finetuning approach that significantly reduces parameter requirements. BOFT encompasses OFT as a special instance, thus offering a unified and generalized orthogonal finetuning paradigm. We further substantiate our approach with comprehensive experiments on the adaptation of large vision transformers, language models, and text-to-image diffusion models to a diverse array of vision and language downstream applications.
\end{abstract}

\section{Introduction}

Prominent models like ChatGPT~\cite{brown2020language,chen2021evaluating} and Stable Diffusion~\cite{rombach2022high} showcase the extraordinary ability of large foundation models to generalize across tasks. This leap in performance is closely tied to a sharp rise in parameter counts—

(\emph{e.g.}, GPT-3 features approximately 175B parameters)—which has made full training of such models increasingly prohibitive for the research community. Advancing the field thus hinges on efficiently repurposing existing foundation models rather than building them anew. In particular, effective and efficient adaptation to downstream tasks is vital. There are three predominant strategies to achieve this: \emph{model finetuning}~\cite{hulora2022,zaken2022bitfit,guo2020parameter,raffel2020exploring,qiu2023controlling,zhang2022adaptive,chavan2023one}, where only select model weights are updated while the original architecture remains intact; \emph{adapter tuning}~\cite{rebuffi2017learning,houlsby2019parameter,pfeiffer2020adapterfusion,lin2020exploring,he2021towards}, which introduces auxiliary trainable components alongside the base model; and \emph{prompt tuning}~\cite{li2021prefix,lester2021power}, which augments the input with tunable prefix tokens. Among these, model finetuning is notable for its simplicity and robust effectiveness, offering the additional benefit of maintaining the model's inference speed by not introducing extra computational overhead.

The central idea underlying model finetuning is to maintain a close resemblance between the pretrained and finetuned models according to specific criteria, thereby safeguarding the knowledge acquired during pretraining. As an example, many current finetuning techniques employ a low learning rate; this heuristic serves to keep modifications between the pretrained and finetuned models minimal. Recognizing the inherent structure within weight matrices, more systematic approaches focus on preserving relational characteristics, such as maintaining the pairwise angular relationships among neurons. This observation has inspired the development of a new finetuning methodology called Orthogonal Finetuning~(OFT)~\cite{qiu2023controlling}. In OFT, neurons within a given layer are updated using a shared orthogonal transformation, guaranteeing the preservation of pairwise angles throughout finetuning. Despite OFT's demonstrated benefits for generalization and optimization when applied to finetuning text-to-image diffusion models~\cite{qiu2023controlling}, it requires a substantial number of trainable parameters, owing to the size of orthogonal matrices involved. To mitigate this, OFT imposes a block-diagonal structure on the orthogonal matrix to decrease parameter count. Nonetheless, this gain in parameter efficiency brings certain drawbacks: the imposed block-diagonal sparsity constrains the orthogonal transformation to fixed, non-overlapping blocks, potentially reducing representational capacity and limiting the matrix’s ability to approximate general linear operators. More critically, there is no established method for selecting an optimal sparsity structure, leaving this as an open challenge.

The central challenge lies in constructing a dense orthogonal matrix in a way that minimizes the number of parameters. At first, this appears daunting, as a dense orthogonal matrix of size $d$ requires $\mathcal{O}(d^2)$ parameters. Instead of following the standard approach, we introduce a strategy that builds a dense orthogonal matrix by multiplying a sequence of sparse, factorized matrices. This methodology is motivated by the idea that one can conserve parameters by incurring additional computational cost. Because the orthogonal matrix is expressed as a product of sparse matrices, this sequence of multiplications must be carried out at every step during training. To clarify this matrix decomposition, we reframe dense orthogonal matrix construction as a problem of information propagation. Viewed from this angle, forming a dense orthogonal matrix via sparse matrix multiplication can be likened to propagating information across a grid-structured network. This perspective enables the creation of various sparse matrix factorizations that keep the number of learnable parameters low, yet retain enough flexibility to produce fully dense matrices.

For parameter-efficient design within this information propagation paradigm, we are inspired by the butterfly patterns found in the Cooley-Tukey fast Fourier transform algorithm~\cite{cooley1965algorithm}, where rapid exchange of information among nodes is possible~\cite{klasing1994broadcasting}. The butterfly network inherent in the Cooley-Tukey algorithm gives rise to a sparse matrix factorization scheme, known as butterfly factorization. Using butterfly factorization, we can generate a $d$-dimensional dense matrix as a product of $\mathcal{O}(\log d)$ sparse factors, where each has only $\mathcal{O}(d)$ non-zero elements. By enforcing orthogonality for each factor, we achieve a compact orthogonal parameterization, requiring only $\mathcal{O}(d\log d)$ parameters—a considerable reduction compared to standard dense representations. Building on this compact structure, we introduce Orthogonal Butterfly (BOFT), a new method for parameter-efficient finetuning. BOFT generalizes OFT, encompassing it as a particular instance, and constitutes a broad framework for orthogonal finetuning. Notably, both block-diagonal and butterfly patterns inject sparsity into orthogonal matrices, decreasing the total number of adjustable parameters. It is worthwhile to juxtapose this approach with the low-rank structure used in LoRA~\cite{hulora2022}; both low-rank and sparse matrices represent principal categories of structured matrices~\cite{candes2011robust} that foster parameter efficiency.

In contrast to the block-diagonal architecture leveraged by OFT to balance between the model's expressiveness and regularization, BOFT employs the butterfly structure to achieve a more gradual interpolation from the full orthogonal group of matrices (which affords greater expressivity) toward the identity matrix (imposing regularity). This structural design permits the identification of a reduced hypothesis space within the orthogonal group, which is beneficial for subsequent learning tasks. Given the butterfly structure's prevalent use in rapid linear transformations—such as the discrete Fourier transform and the discrete cosine transform—we posit that imposing this structure in parameter optimization serves as an implicit inductive bias, likely enhancing generalization and curbing overfitting. Our rationale is anchored in the observation that while the butterfly architecture can naturally reproduce various classical linear transformations, the block-diagonal approach of OFT cannot. The principal contributions of our work are as follows:

\begin{itemize}[leftmargin=*,nosep]
    \item We introduce a new information transmission framework to explore parameter-efficient orthogonal finetuning. In this framework, designing a dense orthogonal matrix with high parameter efficiency is reformulated as solving an information transmission task on a grid-structured graph.
    \item Drawing inspiration from butterfly structures in the Cooley-Tukey algorithm, we put forth Orthogonal Butterfly—a parameter-efficient method for orthogonal finetuning—framed within the information transmission perspective.
    \item We offer several theoretical perspectives elucidating how BOFT achieves substantial reductions in trainable parameters without compromising model expressivity or training stability. Additionally, the matrix factorization inherent to BOFT leads to a novel mechanism for interpolating weights (refer to Figure~\ref{fig:boft_interpolation}).
    \item We pioneer the use of orthogonal finetuning in diverse domains beyond controllable text-to-image generation~\cite{qiu2023controlling}, demonstrating its promise as a broadly applicable model adaptation strategy. To illustrate, we apply BOFT across adaptation tasks in computer vision and natural language processing, consistently surpassing the performance of leading approaches, thereby highlighting its exceptional parameter efficiency and generalization capabilities.
\end{itemize}

\section{Related Work}

\textbf{Parameter-efficient finetuning (PEFT)}. With the rapid scaling of foundation models (\emph{e.g.}, \cite{brown2020language,radford2021learning,rombach2022high,kirillov2023segment}), developing methods to adapt these models to specific tasks while minimizing the number of updated parameters has become an area of active research~\cite{treviso2023efficient,ding2023parameter,lialin2023scaling}. Among the array of PEFT approaches~\cite{houlsby2019parameter,aghajanyan2020intrinsic,hulora2022,edalati2022krona,wang2022adamix,gheini2021cross,zaken2022bitfit,guo2020parameter,sung2021training,ansell2022composable,lester2021power,li2021prefix,vu2022spot,he2021towards,mao2021unipelt,karimi2021compacter,liu2022few,sung2022lst,chen2022parameter,jia2022visual,chen2022adaptformer,zhang2022neural,jie2023fact,lian2022scaling,luo2023towards}, techniques based on reparameterization~\cite{aghajanyan2020intrinsic,hulora2022,edalati2022krona} align most closely with our objectives. LoRA~\cite{hulora2022}, for example, modifies pretrained weights by incorporating the product of two low-rank matrices, leading to strong results on a variety of language tasks. Given that LoRA employs a constant rank for all inserted matrices, subsequent works~\cite{zhang2022adaptive,valipour2022dylora,zhang2023increlora} have proposed adaptively setting the rank at each layer, which leads to better allocation of limited parameters. Owing to its straightforwardness, this type of low-rank reparameterization has seen widespread adoption~\cite{dettmers2023qlora,zi2023delta,chavan2023one}. Motivated by the role of hyperspherical energy in capturing generalization properties~\cite{liu2018learning,liu2021orthogonal}, \cite{qiu2023controlling} introduces orthogonal finetuning (OFT) as a distinctive and successful approach for adapting text-to-image diffusion models. In OFT, the finetuning process involves learning an orthogonal transformation within a given layer, which has been empirically shown to yield better generalization and more stable training than LoRA. Nonetheless, this improvement typically requires more trainable parameters than LoRA. Enhancing the parameter efficiency of OFT therefore remains a meaningful challenge. Furthermore, it remains uncertain if OFT extends effectively to a broader array of adaptation scenarios beyond text-to-image diffusion~\cite{qiu2023controlling}. BOFT addresses this by integrating butterfly factorization, thereby increasing the efficiency of OFT. This allows us to systematically demonstrate the effectiveness of orthogonal finetuning across diverse adaptation tasks.

\textbf{Butterfly structures}. The radix-2 Cooley-Tukey algorithm decomposes the $N$-point discrete Fourier transform into two smaller $\frac{N}{2}$-point transforms at each step, yielding a computational pattern known as the butterfly structure. This pattern is expressible as a sequence of sparse matrix multiplications, often referred to collectively as a butterfly matrix. Applications of butterfly matrices include parameterizing orthogonal transformations to obviate the need for pivoting in Gaussian elimination and to enhance computational speed~\cite{parker1995random}, facilitating stable training in recurrent neural networks~\cite{jing2017tunable}, and advancing kernel approximation techniques~\cite{munkhoeva2018quadrature}. Research by \cite{dao2019learning,dao2019kaleidoscope} has explored the learning of efficient linear operators via butterfly parameterizations. Meanwhile, \cite{chen2022pixelated,dao2022monarch} have harnessed butterfly matrices for more efficient neural network training. Additional utilizations of butterfly structures have been demonstrated in the context of accelerated matrix-vector operations~\cite{michielssen1996multilevel,o2010algorithm}, in compressing data-sparse matrices~\cite{li2015butterfly}, and for improved network transmission~\cite{klasing1994broadcasting,li2003linear,rajasingh2016transmission}. Departing from these lines of inquiry, our investigation centers on leveraging butterfly structures specifically to boost the parameter efficiency of OFT.

\section{An Information Transmission View on Orthogonal Finetuning}

We begin by outlining fundamental concepts related to OFT. The OFT approach adapts a pretrained weight matrix by expressing the updated matrix as the product of a trainable orthogonal matrix and the fixed pretrained matrix. In contrast to LoRA, which relies on adding a low-rank matrix to the weights, OFT applies a multiplicative orthogonal matrix for weight updates. Both methods aim for parameter efficiency: LoRA achieves this through its low-rank design, whereas the original OFT adopts a block-diagonal structure for the orthogonal matrix~\cite{qiu2023controlling} (where smaller block sizes correspond to greater parameter savings). Figure~\ref{fig:comp_lora} offers an intuitive visualization of these differences. The rationale for utilizing orthogonal transformations during finetuning lies in their ability to maintain the angles between neuron vectors~\cite{liu2018learning,liu2021orthogonal,qiu2023controlling}, thereby helping to retain semantic information acquired during pretraining. Specifically, OFT learns an orthogonal matrix $\bm{R}\in\mathbb{R}^{d\times d}$ in conjunction with a pretrained linear layer $\bm{W}^0\in\mathbb{R}^{d\times n}$, replacing the usual forward computation $\bm{z}=(\bm{W}^0)^\top\bm{x}$ with $\bm{z}=(\bm{R}\bm{W}^0)^\top\bm{x}$, where $\bm{x}\in\mathbb{R}^{d}$ and $\bm{z}\in\mathbb{R}^{n}$ denote the input and output vectors, respectively. To ensure zero initialization, OFT sets $\bm{R}$ to the identity matrix at the outset. Maintaining the orthogonality of $\bm{R}$ throughout finetuning is handled by leveraging Cayley parameterization as described in~\cite{qiu2023controlling,liu2021orthogonal}: $\bm{R}=(\bm{I}+\bm{Q})(\bm{I}-\bm{Q})^{-1}$, with $\bm{Q}$ being a skew-symmetric matrix ($\bm{Q}=-\bm{Q}^\top$). To promote parameter efficiency, OFT further enforces a block-diagonal architecture for $\bm{R}$ by setting $\bm{R}=\text{diag}(\bm{R}_1,\bm{R}_2,\cdots,\bm{R}_r)$, where each $\bm{R}_i\in\mathcal{R}^{b\times b}$ is a small orthogonal block and $br=d$. While this block-diagonal design reduces parameter count, it entails that the neuron's dimensions (the columns of $\bm{W}^0$) are partitioned into $r$ groups, each being transformed independently by its corresponding orthogonal block. Although this blockwise approach has shown strong empirical performance, partitioning neuron dimensions purely by index lacks a meaningful basis, thereby motivating the consideration of fully dense orthogonal matrices. This raises a key research question: \emph{Is it possible to devise a dense orthogonal matrix that retains the parameter-efficiency characteristic?}

To tackle this issue, we suggest constructing a dense orthogonal matrix by multiplying several sparse orthogonal matrices together. To present a cohesive and accessible way to analyze the sparsity structures in orthogonal matrix decompositions, we reinterpret the formation of a dense orthogonal matrix within OFT as analogous to a process of information transmission. Concretely, forming a dense matrix $\bm{R}\in\mathbb{R}^{d\times d}$ through the multiplication of $m$ square matrices as $\bm{R}=\bm{B}_m\bm{B}_{m-1}\cdots\bm{B}_1$ can be conceptualized as passing information across a grid comprised of $d\times (m+1)$ nodes, as depicted in Figure~\ref{fig:info_trans}. This perspective is inspired by the observation that a dense $d$-dimensional square matrix captures full connectivity between one set of $d$ nodes and another. In the case of $\bm{R}$, an entry $R_{ij}=0$ means there is no information transfer from the $j$-th input node to the $i$-th output node, while a nonzero value of $R_{ij}$ allows transmission. Therefore, expressing the dense matrix $\bm{R}$ as the product $\bm{B}_m\bm{B}_{m-1}\cdots\bm{B}_1$ can be viewed as sequentially propagating information according to the connectivity encoded in each factor $\bm{B}_i$. The path of information starts with the connections specified by $\bm{B}_1$ and progresses through to those defined by $\bm{B}_n$. As shown in Figure~\ref{fig:info_trans}, an explicit example is given for $\bm{R}=\bm{B}_5\bm{B}_4\bm{B}_3\bm{B}_2\bm{B}_1$, with their respective sparsity structures and the corresponding induced graph. This graph results from unfolding the matrix product, where each $\bm{B}_i$ functions as a connectivity map between nodes on adjacent layers: specifically, an entry $(j_1,j_2)$ in $\bm{B}_i$ specifies whether a directed edge runs from the $j_2$-th node at level $i$ to the $j_1$-th node at level $i+1$ (a zero indicating absence of such a link). For the product $\bm{B}_5\bm{B}_4\bm{B}_3\bm{B}_2\bm{B}_1$ to be dense, every node at the initial level must have a communication path to each node in the sixth level. When only $\bm{R}=\bm{B}_4\bm{B}_3\bm{B}_2\bm{B}_1$ is considered—connecting nodes at the first and fifth levels—we can observe, for instance, that node $1$ is unable to reach node $3$. This reveals a zero at position $(3,1)$ in the sparsity pattern of the product $\bm{B}_4\bm{B}_3\bm{B}_2\bm{B}_1$. Analogous relationships between the underlying graph and the zero structure can be found when considering $\bm{R}=\bm{B}_3\bm{B}_2\bm{B}_1$ or $\bm{R}=\bm{B}_2\bm{B}_1$. In general terms, for $\bm{R}\in\mathbb{R}^{d\times d}$ to be dense, the sequence of matrices $\bm{B}_m,\cdots,\bm{B}_1$ must establish a configuration of directed edges across a $d\times (m+1)$ grid, where each edge joins nodes from adjacent columns (levels). This ensures that information from any node in the first layer can traverse to all nodes at the $(m+1)$-th layer.

Interpreting the sparsity of factorization matrices in OFT through the lens of information transmission offers valuable insights. As depicted in Figure~\ref{fig:blk_diag}, $\bm{R}$ in standard OFT possesses a block-diagonal pattern. While such a structure leads to a reduction in the trainable parameters, it cannot yield a fully dense $\bm{R}$. Our objective, instead, is to approximate a dense orthogonal matrix by composing it from $m$ sparse orthogonal components, all the while minimizing the number of free trainable parameters. The information transmission perspective highlights two central requirements: \emph{(i)} \textbf{dense connectivity}, ensuring that every node at the initial layer is connected to every node in the final layer via at least one route; and \emph{(ii)} \textbf{minimum free edges}, where the cumulative number of edges should be minimized without violating the orthogonality constraint. Orthogonality itself imposes intricate restrictions on how edges may connect neighboring layers. Notably, each $\bm{B}_i$ must remain full-rank (essential for orthogonality), which demands exactly $d$ edges to set up a bijective mapping from the $i$-th to $(i+1)$-th layer, requiring at minimum $d$ connections between adjacent levels (as seen, for instance, in Figure~\ref{fig:info_trans}). Since these $d$ mandatory edges, tied to non-trainable entries (for instance, in a $d\times d$ orthogonal matrix these must be $\pm 1$), do not contribute to model flexibility, they are not included in the count of free edges. Because orthogonal matrices can be specified with fewer parameters than dense matrices, enforcing orthogonality introduces complex dependencies among the free edges. For example, in the case of a $2\times 2$ orthogonal matrix, setting the $(1,1)$ position to zero necessitates the $(2,2)$ entry also being zero; in other words, one absent connection can determine the absence of another. Thus, the orthogonality condition may alter the viable set of connections, at times requiring additional edges or making some redundant. By examining the structure of non-zero elements in the resulting orthogonal matrix, this information transmission viewpoint proves especially pertinent for OFT, since only the locations of the non-zero, trainable entries in $\bm{R}$ matter—their precise values are secondary.

A straightforward approach to establishing dense connections between two levels requires $\mathcal{O}(d^2)$ edges—essentially, this corresponds to implementing a single dense orthogonal matrix—which results in $d^2-d$ distinct edges (for instance, when $d=4$, there are 12 edges). As illustrated in Figure~\ref{fig:info_trans}, a concrete example of a permissible matrix factorization uses only $10$ edges overall, which is fewer than those required by a dense orthogonal matrix. This analytical setup allows us to assess the parameter-efficiency of OFT graphically and facilitates the construction of valid matrix factorizations under this perspective. The structure draws upon a notable topology from the Cooley-Tukey algorithm, namely butterfly graphs~\cite{cooley1965algorithm}, which enable dense connectivity between $d$ sources and $d$ destinations using only $\mathcal{O}(d\log d)$ edges. For illustration, while the network in Figure~\ref{fig:info_trans} leverages $10$ edges for full connectivity, a butterfly topology can achieve the same with merely $8$ edges. In the following, we describe how adopting the butterfly structure yields gains in parameter efficiency.

\section{Orthogonal Parameterization by Butterfly Factorization}

Initially introduced within the Cooley-Tukey algorithm for efficiently computing the fast Fourier transform, the butterfly structure ensures that a local adjustment in the frequency space is reflected globally in the spatial domain—closely analogous to our setting, where every node in the initial layer is required to communicate with all nodes in the terminal layer. The butterfly architecture also underpins prevalent network topologies in computer systems~\cite{solihin2015fundamentals,li2011linear}, valued for its capability to facilitate high-throughput information exchange. Letting $k\geq 2$ denote a power of $2$, we define the \emph{butterfly factor} $\bm{B}^F(k)$ by

\begin{equation}\label{eq:but_factor}
\begin{aligned}
    \bm{B}^F(k)=\begin{bmatrix}
    \text{diag}(\bm{d}_1) & \text{diag}(\bm{d}_2) \\
    \text{diag}(\bm{d}_3) & \text{diag}(\bm{d}_4)
    \end{bmatrix}\in\mathbb{R}^{k\times k},
\end{aligned}
\end{equation}

where each $\bm{d}_i\in\mathbb{R}^{\frac{k}{2}}$ for $i=1,\dots,4$ denotes a vector. If $d=2^N$, the $d$-dimensional \emph{butterfly matrix} $\bm{B}(d)\in\mathbb{R}^{d\times d}$ is constructed recursively by

\begin{equation}
\begin{aligned}
    \bm{B}(d)=\tilde{\bm{B}}(d,d)\!\cdot\!\begin{bmatrix}
    \bm{B}_1(\frac{d}{2})\!\!\!\!\! & \bm{0} \\
    \bm{0}\!\!\!\!\! &\bm{B}_2(\frac{d}{2})
    \end{bmatrix}= \tilde{\bm{B}}(d,d)\tilde{\bm{B}}(d,\frac{d}{2})\cdots\tilde{\bm{B}}(d,2),
\end{aligned}
\end{equation}

where $\bm{B}_1(\frac{d}{2})$ and $\bm{B}_2(\frac{d}{2})$ denote two butterfly matrices, each of dimension $\frac{d}{2}$. We introduce the \emph{butterfly component} as $\thickmuskip=2mu \medmuskip=2mu \tilde{\bm{B}}(d,k)=\text{diag}(\bm{B}^F_1(k),\cdots,\bm{B}^F_{d/k}(k))$, a block-diagonal matrix of dimension $d \times d$ composed of blocks of size $k$. Here, $\bm{B}^F_i(k)$ (for each $i$) are the butterfly factors described in Equation~\ref{eq:but_factor}. With these definitions, we proceed to construct an orthogonal matrix using the butterfly matrix framework. The only requirement is that each factor $\tilde{\bm{B}}(d,k)$ comprising the butterfly matrix $\bm{B}(d)$ is orthogonal. Starting with the simplest case, we examine the block-diagonal matrix $\tilde{\bm{B}}(d,2)$ with block size $2$; here, orthogonality is trivially ensured by parameterizing each block using the Cayley transform or a $2$-dimensional rotation, following the method in \cite{liu2021orthogonal,qiu2023controlling}. The pattern of non-zeros in any butterfly component can be obtained by permuting the non-zero configuration of $\tilde{\bm{B}}(d,2)$; thus, every butterfly component can similarly be made orthogonal through analogous parameterization. This approach yields a highly efficient orthogonal matrix representation constructed from multiple $2\times 2$ orthogonal submatrices. Extending the construction as in \cite{chen2022pixelated}, we define a block butterfly component $\tilde{\bm{B}}^b(d,k)$ where each element in $\bm{d}_i$ is replaced by a $b\times b$ matrix. To enforce orthogonality of $\tilde{\bm{B}}^b(d,2)$, each $2b\times 2b$ block matrix is individually parameterized to be orthogonal. For $\tilde{\bm{B}}^b(d,k)$ with $k>2$, the structure of non-zero entries results from a block-wise permutation of the configuration found in $\tilde{\bm{B}}^b(d,2)$, allowing these larger components to also be rendered orthogonal using a similar strategy. Integrating these components, the BOFT forward pass is given by

\begin{equation}
    \begin{aligned}\nonumber
    \bm{z}=\big(\bm{R}(m,b)\cdot\bm{W}^0\big)^\top\bm{x},~~~~\text{s.t.}~\bigg{\{}\bm{R}(m,b)=\prod_{i=1}^m \tilde{\bm{B}}^b_{(d,i)}~~\&~~\underbrace{\big(\tilde{\bm{B}}^b_{(d,j)}\big)^\top\tilde{\bm{B}}^b_{(d,j)}=\tilde{\bm{B}}^b_{(d,j)}\big(\tilde{\bm{B}}^b_{(d,j)}\big)^\top\!=\bm{I}_d}_{\forall j\in [1, m]}\bigg{\}},
    \end{aligned}
\end{equation}

Here, $\tilde{\bm{B}}^b(d,2^{m - i+1})$ is simplified and written as $\tilde{\bm{B}}^b_{(d,i)}$, and $\bm{I}_d$ designates the $d \times d$ identity matrix. The orthogonal matrix $\bm{R}(m,b)\in\mathbb{R}^{d\times d}$ is constructed by multiplying several orthogonal butterfly matrices together. For clarity, we use BOFT with $\bm{R}(m,\frac{b}{2})$ as shorthand for BOFT($m$,$b$), assuming $b\geq 2$. If $m=1$, BOFT($1$,$b$) simplifies to the block-diagonal OFT~\cite{qiu2023controlling} with block size $b$. Similarly, setting $b=d$ in BOFT($1$,$d$) results in the conventional OFT~\cite{qiu2023controlling} with a fully unconstrained orthogonal matrix. In the case of BOFT($\log \frac{2d}{b}$,$b$), a block butterfly matrix $\bm{B}^{\frac{b}{2}}(d)$ serves as $\bm{R}$, leading to a dense orthogonal transformation. In the general case, BOFT($m$,$b$) utilizes $\frac{1}{2}(b-1)dm$ effective trainable parameters for finetuning a linear layer of size $d\times n$. If one sets $m = \log d$ and $b = 2$, corresponding to the butterfly matrix, BOFT then requires only $\mathcal{O}(d\log d)$ parameters. In contrast, using a full dense orthogonal matrix in OFT leads to $\mathcal{O}(d^2)$ parameters, while a block-diagonal OFT with $r$ blocks entails $\mathcal{O}(bd)$ parameters. Consequently, the standard OFT needs $b=d$ for a dense orthogonal matrix, but BOFT can create such matrices for any block size $b$.

\textbf{Identity initialization for BOFT}. Common finetuning procedures often initialize with weights from the original pretrained network to minimize deviation during adaptation; for example, LoRA uses zero initialization in its low-rank modules. In BOFT, every butterfly matrix is initialized to be the identity, meaning the skew-symmetric matrix in the Cayley transform is set to zero.

\textbf{Multiplicative Dropout for BOFT}. LoRA~\cite{hulora2022} integrates Dropout into its low-rank updates as a regularization strategy. While traditional Dropout~\cite{srivastava2014dropout} integrates seamlessly with LoRA, it is not naturally compatible with BOFT’s multiplicative update scheme. To tackle this challenge, we develop a version of multiplicative Dropout specific to BOFT. Since $\bm{R}(m,b)$ consists of $m$ orthogonal butterfly matrices, each capable of being restructured into $2b\times 2b$ block-diagonal forms, multiplicative Dropout is implemented by randomly selecting $p_1\%$ of the butterfly components and $p_2\%$ of their diagonal blocks and replacing these selected submatrices with identity blocks.

\section{Intriguing Insights and Discussions}

\textbf{Expressivity of BOFT}. The butterfly-based structure combined with permutations is capable of exactly reconstructing several classical fast linear transforms~\cite{dao2019learning,dao2019kaleidoscope} (such as the fast Fourier transform and the Hadamard transform). Nevertheless, the extent to which our orthogonal butterfly matrix can approximate arbitrary orthogonal matrices is still an open question. To investigate this, we carried out simulations aiming to approximate a randomly sampled dense orthogonal matrix~\cite{anderson1987generation} of dimension $512\times 512$. Results presented in Figure~\ref{fig:para_eff} represent averages over 10 random seeds. Here, the y-axis reflects the approximation error, while the x-axis captures the count of effective trainable parameters. Within each color, curves denote BOFT configurations with a fixed block size; the leftmost marker for each curve corresponds to the BOFT($1$,$b$) case (that is, the basic block-diagonal OFT with block size $b$). Across various configurations, BOFT typically achieves superior parameter efficiency relative to OFT. As an illustration, BOFT($9$,$2$) demonstrates higher expressivity than BOFT($1$,$16$) despite utilizing far fewer parameters. Generally, adopting smaller $b$ and larger $m$ within BOFT yields greater parameter efficiency; for instance, BOFT($6$,$4$) employs significantly fewer parameters yet attains approximation error similar to BOFT($2$,$16$). In summary, the butterfly matrix encodes a more structured subset within the orthogonal group—when contrasted with a block-diagonal form—which endows BOFT with provably higher expressivity than OFT when both use the same block size.

\begin{theorem}[Expressivity of BOFT]\label{thm:exp_boft}
    BOFT is strictly more expressive than OFT with equivalent block size. Specifically, in order for a butterfly matrix to approximate every orthogonal matrix of size $d$, it suffices to multiply butterfly matrices as $\bm{B}_{d-1,1}(d)\bm{B}^\top_{d-1,2}(d)\cdots\bm{B}_{1,1}(d)\bm{B}^\top_{1,2}(d)$, where all $\bm{B}_{i,j}(d)$ represent butterfly matrices for any $i$ and $j$.
\end{theorem}

Theorem~\ref{thm:exp_boft} motivates a natural generalization for BOFT. Here, the resulting orthogonal matrix is given as $\thickmuskip=2mu \medmuskip=2mu \bm{R}^G(m_1,b_1,m_2,b_2,l)=\bm{R}_{l,1}(m_1,b_1)\bm{R}_{l,2}^T(m_2,b_2)\cdots\bm{R}_{1,1}(m_1,b_1)\bm{R}_{1,2}^T(m_2,b_2)$, where $\bm{R}_{i,j}^T(m,b)$ denotes the orthogonal matrices utilized within BOFT. By selecting $m_1=m_2=\log d$, $b_1=b_2=2$, and $l=d-1$, the resulting $\bm{R}^G(m_1,b_1,m_2,b_2,l)$ spans the full orthogonal group, an arrangement also recognized as the kaleidoscope hierarchy~\cite{dao2019kaleidoscope}. Notably, although such increased expressivity expands the representational scope, it does not invariably yield better empirical results during finetuning. For example, full finetuning—despite its maximal expressivity—often achieves less-than-desired outcomes. Therefore, balancing expressivity and regularity is essential for ensuring the generalization capability of model finetuning. By leveraging the structured priors of the butterfly construction, BOFT expands the space of possible finetuning parameters and thus enables a more advantageous trade-off between expressivity and regularization.

\textbf{Spectral properties}. Orthogonal finetuning tends to achieve superior spectral characteristics compared to LoRA, primarily because it exactly maintains the spectral norm of the pretrained parameter matrix $\bm{W}^0$. To see this, consider the singular value decomposition: $\bm{W}^0=\bm{U}\bm{\Sigma}\bm{V}^\top$, where $\bm{U}$ and $\bm{V}$ are orthogonal matrices, and $\bm{\Sigma}$ is diagonal with singular values. Both OFT and BOFT involve multiplying an orthogonal matrix $\bm{R}$ on the left, yielding updated weights of the form $\bm{R}\bm{U}\bm{\Sigma}\bm{V}^\top$. This operation leaves the maximum singular value—equivalently, the spectral norm—of $\bm{W}^0$ unchanged. Preserving this property has been empirically demonstrated to enhance both the stability of training and model generalization~\cite{miyato2018spectral,yoshida2017spectral}. Additional mathematical insights regarding this preservation are presented in Appendix~\ref{app:sn_property}.

\textbf{Orthogonal finetuning as learning bilinear similarity}. BOFT may also be interpreted through the lens of bilinear similarity, as it can be represented by $\bm{w}^0_i\bm{R}\bm{x}$, where $\bm{w}^0_i$ denotes the $i$-th column (i.e., neuron) from $\bm{W}^0$. Thus, BOFT essentially involves learning the bilinear similarity matrix $\bm{R}$ under the constraint that $\bm{R}$ is orthogonal, establishing a direct link to the field of distance metric learning~\cite{xing2002distance} and the theory of bilinear forms~\cite{roman2005advanced}.

\textbf{Inductive bias and generalization in BOFT}. In the context of BOFT, the matrix $\bm{R}(m,b)$ typically belongs to a structured subclass of the orthogonal group, which restricts the function class and introduces a particular inductive bias. We propose that the structured bias from butterfly factorization contributes positively to generalization. This arises because the butterfly factorization framework shares structural motifs with various standard linear transforms~\cite{dao2019kaleidoscope}, including the discrete Fourier, sine/cosine, and Hadamard transforms. Furthermore, the use of sparse matrix factorization in BOFT can introduce additional implicit inductive biases~\cite{gunasekar2017implicit,li2018algorithmic,liu2019neural}.

\textbf{Comparison to butterfly-based sparse training}. Numerous prior studies~\cite{dao2019learning,dao2019kaleidoscope,chen2022pixelated,dao2022monarch} have investigated sparse training methods that utilize butterfly parameterizations. These works generally employ butterfly structures to reparameterize the weight matrices and proceed to train the networks from scratch. In contrast, \cite{dao2022monarch} explores adapting pretrained models by initially projecting their weights onto modified butterfly matrices, followed by optimizing the resulting components for specific downstream applications. The BOFT approach, on the other hand, introduces a distinct finetuning methodology, leveraging shared transformation matrices across layers.

\input{rebuttal-results}

\section{Concluding Remarks and Limitations}

In this work, we introduce Orthogonal Butterfly, a general-purpose, parameter-efficient finetuning technique for foundation models built upon the butterfly architecture. Our main conceptual advance lies in expressing a full orthogonal matrix as the product of several sparse orthogonal matrices, enhancing parameter efficiency. To facilitate practical matrix decompositions, we develop a graphical information transmission perspective. Within this paradigm, we observe that the butterfly design provides an effective mechanism for sparse orthogonal factorization. BOFT is empirically validated across a variety of tasks, including the adaptation of large language models, vision models, and text-to-image generators, establishing its advantage as a flexible and efficient finetuning solution.

Nevertheless, BOFT exhibits some limitations. In particular, due to the composition of multiple orthogonal matrices to construct the final transformation, training incurs a moderate overhead compared to OFT. Mitigating BOFT's training runtime cost remains an unresolved challenge. Encouragingly, after finetuning, the orthogonal matrices derived via BOFT can be merged into the original model without imposing additional latency during inference. Furthermore, it is still unclear if butterfly networks represent the optimal architecture for information propagation. Our proposed information transmission viewpoint also opens avenues for incorporating insights from computer networking literature, where topological efficiency for information exchange is a central concern. We anticipate that adopting more advanced network topologies could further enhance the construction of dense orthogonal matrices.